\par\Needspace{7\baselineskip}\bigskip
\begin{center}
\SourceLeaf{442}
\pdfbookmark[2]{Gauss an Hansen, Göttingen, 11. März 1843}{pallas-brief-18430311}
\textsc{Gauss} an \textsc{Hansen}, Göttingen, 11. März 1843.
\end{center}
\noindent\hspace*{\parindent}Von einer Woche zur anderen ist mein Dank für die gefällige Übersendung des Berliner Monatsberichts (der um dieselbe Zeit auch auf gewöhnlichem Wege mir zu Gesicht kam) verschoben, weil ich hoffte, einige Zeit zu gewinnen, in den Gegenstand etwas weiter in meiner Antwort eingehen zu können. Leider ist diese Hoffnung getäuscht, und selbst in den bevorstehenden Ferien, für welche sich schon im Voraus so viele Rückstände und neue Abhaltungen gesammelt haben, darf ich mir kaum Hoffnung für einige freie Muße machen.
\par Ich beschränke mich daher darauf, meine Freude darüber auszusprechen, dass Sie bei den Perturbationsrechnungen auf ähnliche Art verfahren, wie ich schon vor mehr als 30 Jahren bei meinen weitumfassenden Rechnungen über die Pallasstörungen zu Werke gegangen bin, insofern Sie den Gebrauch von Reihen nach den Exzentrizitäten und Neigungen ganz kassieren. Freilich haben Sie für die Kometenstörungen auch ganz besondere noch andere Untersuchungen nötig gehabt, zu denen für die Pallas keine Veranlassung sich fand, und überhaupt in vielen anderen Beziehungen abweichende Wege eingeschlagen. Bei den Störungen der Pallas durch Jupiter brauchte ich die Methode der variablen Elemente, und zwar mit Vorbedacht, denn obgleich man so 6 Elemente zu behandeln hat, während bei dem anderen Verfahren nur halb so viele sind, so habe ich doch für den praktischen Gebrauch jenes vorgezogen. Man braucht die Rechnung für jedes Jahr nur einmal (für Einen Tag) zu machen, wozu ich eine besondere Hilfstafel konstruierte, vermittelst welcher in vergleichsweise sehr kurzer Zeit die Rechnung absolviert werden kann, obgleich zusammen 801 Gleichungen (1602, wenn die Sinus- und Kosinusglieder desselben Arguments getrennt gezählt werden. Es sind vollständig alle, deren Koeffizient über \(0''{,}1\) geht) berücksichtigt worden\footnote{Es ist wohl unnötig zu bemerken, dass die Arbeit, von der ich jetzt spreche, \textls[180]{ganz} verschieden ist von der gleichfalls von mir für den ganzen Zeitraum von 1802 bis etwa 1818 oder 1820 fortgeführten Rechnung durch Quadratur. Diese nenne ich spezielle, jene generelle Rechnung und letztere hat in jener eine bei so ausgedehnten Rechnungen höchst notwendige Kontrolle gefunden.}. Durch sehr einfache Mittel kann man dann das Resultat für die ganze \SourceLeaf{443}Beobachtungssaison ausreichend machen, ohne der Schärfe etwas zu vergeben. Auch die Störung durch Saturn wurde berechnet und die durch Mars nach einer wesentlich verschiedenen Methode angefangen, aber nicht vollendet. Andere immer weiter sich verzweigende Geschäfte haben mir später gar nicht erlaubt, auf jene Arbeiten wieder zurückzukommen, und es steht dahin, ob in meinem Alter ich Muße und Lust haben werde, mich wieder in die Sachen hineinzuarbeiten, da noch so viele andere Dinge sind, die ich ebenso ungern untergehen lassen möchte. Sie sind sehr glücklich, dass Sie in einer äußeren Lage sind, wo Sie Ihre Zeit nicht zu versplittern brauchen, und es wird mir jedenfalls zur Beruhigung gereichen, dass dieser Zweig der Astronomie bei Ihnen in den besten Händen ist. Ob ich bei den Marsstörungen etwas mit Ihrer mir angezeigten Integrationsart Zusammenhängendes gebraucht habe, kann ich jetzt, wo mir die Sachen seit fast 25 Jahren entfremdet sind, nicht bestimmt ermitteln, möchte es aber fast bezweifeln.\quad\mbox{. . . . .}
\par
